\section{Problem Formulation}
\label{sec:formulation}

\begin{algorithm}[htbp]
\small
\caption{Walk-Forward Evaluation}
\label{alg:walkforward}

\KwIn{
OHLCV series $\mathbf{X}$,
alpha registry $\mathcal{A}$,
historical sentiment cache $\mathcal{D}$,
window size $W$,
lookahead $L$,
step size $S$
}

\KwOut{
Directional predictions $\hat{y}_{e,L}$,
accuracy, Alpha Lift, and compounded account return
}

Initialize results list $\mathcal{R} \leftarrow \emptyset$

\ForEach{test point $e$ in the walk-forward schedule}{

    Extract the exact analysis window $\mathbf{X}_e=\mathbf{X}_{[e-W,e)}$

    Compute technical indicators using TA-Lib

    Evaluate all candidate alphas $\alpha^{(k)} \in \mathcal{A}$

    Remove candidates with no fitted transform, fewer than 20 finite
    evaluation observations, or non-finite ranking metrics to form $\mathcal{V}$

    Min--max normalise each metric over $\mathcal{V}$ and rank using:

    \[
    S_k =
    0.40\mathcal{N}_{\mathcal{V}}(|IC_k|)
    + 0.35\mathcal{N}_{\mathcal{V}}(Acc_k)
    + 0.25\mathcal{N}_{\mathcal{V}}(Sharpe_k)
    \]

    Select top-5 alpha factors $\mathcal{A}^{*}$

    Retrieve historical CafeF sentiment at the decision cutoff $d=e-1$

    Generate candlestick and trendline chart images from the same exact
    $\mathbf{X}_e$ window

    Run the Indicator, Alpha, Pattern, and Trend Agents

    Pass all reports to Decision Agent

    Obtain binary prediction
    $\hat{y}_{e,L} \in \{\texttt{LONG}, \texttt{SHORT}\}$

    Compare prediction with realised return

    Store metrics in $\mathcal{R}$
}

Compute directional accuracy and Alpha Lift

Compute compounded account return under Eq.~\eqref{eq:account_return}

\Return{$\mathcal{R}$}

\end{algorithm}
\subsection{Market Model and Data}

Let $\mathcal{S} = \{s_1, s_2, \ldots, s_N\}$ denote the universe of
Vietnamese equities listed on HOSE, HNX, or UPCoM.
For each asset $s \in \mathcal{S}$, the market is observed as a
discrete-time OHLCV series:
\begin{equation}
    \mathbf{x}_u = \bigl(O_u,\, H_u,\, L_u,\, C_u,\, V_u\bigr)
    \;\in\; \mathbb{R}^5_{+},
    \quad u = 0, 1, \ldots, T-1,
\end{equation}
where $O_u$, $H_u$, $L_u$, $C_u$, $V_u$ denote the open, high, low,
closing price, and trading volume at zero-based bar position $u$ respectively.
When volume observations are missing for more than 50\% of bars, the
intra-bar range $H_u - L_u$ is substituted as a liquidity proxy.

We use zero-based row positions and an exclusive window-end boundary.  For a
test point with boundary $e$, the analysis window of size $W$ is exactly:
\begin{equation}
    \mathbf{X}_{e} = \bigl(\mathbf{x}_{e-W},\,
                           \mathbf{x}_{e-W+1},\, \ldots,\,
                           \mathbf{x}_{e-1}\bigr)
    \;\in\; \mathbb{R}^{W \times 5}.
    \qquad\text{(equivalently, }\mathbf{X}_{e}=\mathbf{X}_{[e-W,e)}\text{)}
\end{equation}
Thus $e$ is not itself a visible bar: the final visible (decision) bar is
$d=e-1$.  The next bar, $e$, is reserved for the earliest simulated entry,
and the immutable label target is $h=e-1+L$.  A candidate economic exit may
use that close only if the separate execution contract accepts it. This mapping is used in the
engine, alpha labels, persisted provenance, and the experiments:
\begin{equation}
    \boxed{\quad
    \text{window}=[e-W,e),\qquad d=e-1,\qquad
    \text{entry}=e,\qquad h=e-1+L.\quad}
    \label{eq:time_index_contract}
\end{equation}
The mapping is positional rather than calendar arithmetic; timestamps are
carried alongside every index, so a weekend or holiday gap does not create a
synthetic bar.  The first complete window is $e=W$, and it is evaluable only
when the target is present, i.e., $T\geq W+L$.  For $L=1$, the target is the
entry bar ($h=e$), while for $L=3$ it is two positions after entry
($h=e+2$).  Any boundary whose target is outside the available frame is
rejected rather than shortened.

\subsection{Directional Prediction Task}

The core task predicts whether a candidate LONG cycle is profitable after
the configured transaction costs. With $d=e-1$ and $h=e-1+L$, define
\begin{equation}
    r^{\mathrm{net}}_{e,L}
    =\frac{C_h}{O_e}
      \frac{(1-c_{\mathrm{slip}})(1-c_{\mathrm{fee}})}
           {(1+c_{\mathrm{slip}})(1+c_{\mathrm{fee}})}-1,
    \qquad
    y_{e,L}=
    \begin{cases}
        +1, & r^{\mathrm{net}}_{e,L}>0,\\
        -1, & r^{\mathrm{net}}_{e,L}\leq 0.
    \end{cases}
    \label{eq:label}
\end{equation}
Here $O_e$ is the next-bar entry open and $C_h$ is the target close. The same
net Open-to-Close return defines the directional label, the alpha-selection
target, and the realised return of an executed LONG. Thus an opening gap and
both transaction legs are represented consistently. The reported study uses
daily bars with $L=3$; intraday execution is outside its empirical scope.

A trading system $\mathcal{F}$ maps each analysis window to a binary
decision:
\begin{equation}
    \hat{y}_{e,L} \;=\; \mathcal{F}\!\bigl(\mathbf{X}_e\bigr)
    \;\in\; \{-1,\, +1\},
\end{equation}
where $+1$ corresponds to a \texttt{LONG} signal and $-1$ to
\texttt{SHORT}.

\subsection{Alpha Factor Formalisation and Adaptation}

An \emph{alpha factor} $\alpha^{(k)}$ is a measurable function that maps
a feature frame derived from $\mathbf{X}_e$ to a scalar signal:
\begin{equation}
    \alpha^{(k)} : \mathbb{R}^{T_{\mathrm{hist}} \times d}
    \;\longrightarrow\; \mathbb{R},
\end{equation}
where $d$ is the number of engineered features and
$T_{\mathrm{hist}}$ is the historical window used for computation.

The executable registry contains 80 adaptations drawn from the WorldQuant-101 catalogue, not all 101 source expressions, plus five proprietary formulas. The source expressions were designed primarily for cross-sectional ranking across a broad, liquid equity universe. To apply the registered subset to a single Vietnamese stock, we use a \emph{semantic translation layer} that replaces cross-sectional operators with rolling time-series surrogates. This is an engineering heuristic rather than an invariance-preserving mathematical transformation.

\paragraph{Operator Localisation.}
The primary transformation involves the cross-sectional $\operatorname{rank}(x)$ operator, which denotes the relative percentile of a stock against its peers. A cross-sectional rank is localised to a 20-bar time-series percentile rank unless the source expression already specifies a time-series rank with its own window. Let $\mathcal{H}_{u,w}$ contain the available preceding observations in that rolling window and let $n_{u,w}=|\mathcal{H}_{u,w}|+1$. The implementation computes
\begin{equation}
    \operatorname{ts\_rank}(x_u,w)
    \;=\;
    \frac{1}{\max(n_{u,w}-1,1)}
    \sum_{x_v\in\mathcal{H}_{u,w}} \mathbf{1}[x_u>x_v].
    \label{eq:ts_rank_localised}
\end{equation}
Thus ties are not counted as strictly lower values, and an initial partial window is permitted once the operator-specific minimum is met. Other explicit time-series lookbacks are retained from each registered expression and can be substantially longer than 20 bars. Average Daily Volume ($\mathrm{ADV}_{n}$) is the asset's own $n$-day rolling mean volume, and VWAP is approximated by the daily typical price $(H_u+L_u+C_u)/3$. Functions $\delta(x,d)$ and $\operatorname{delay}(x,d)$ are implemented as $d$-period differences and lags. These substitutions remove peer-relative information, so neither the economic interpretation nor the stochastic properties of the source cross-sectional alpha are presumed to survive unchanged.

\paragraph{Consistent Normalisation.}
To ensure comparability across candidates with different scales, each adapted raw alpha value is normalised to $[-1,1]$ via:
\begin{equation}
    \tilde{\alpha}^{(k)}_{u,e}
    \;=\;
    \tanh\!\left(
        \frac{\alpha^{(k)}_u - \mu_{k,e}}{\sigma_{k,e} + \varepsilon}
    \right),
    \label{eq:alpha_norm}
\end{equation}
where $\mu_{k,e}$ and $\sigma_{k,e}$ are computed over the finite values in
the cutoff-safe calibration snapshot available at decision boundary $e$, and
$\varepsilon > 0$ prevents division by zero. They are recomputed at the next
cutoff; they are not estimated from the entry or target bars.
The resulting signal $\tilde{\alpha}^{(k)}_u \in (-1, 1)$ is treated as
a continuous directional score: positive values favour \texttt{LONG} and
negative values favour \texttt{SHORT}.

\subsection{Walk-Forward Alpha Selection}

Let $\mathcal{A} = \bigl\{\alpha^{(1)}, \alpha^{(2)}, \ldots,
\alpha^{(M)}\bigr\}$ denote the registry of $M = 85$ candidate alpha
functions. Each point-in-time history ends at the decision cutoff and is
therefore disjoint from the current entry and target bars. Given the net
Open-to-Close execution-return series $r^{\mathrm{net},(L)}_u$ defined as in
Eq.~\eqref{eq:label}, the following diagnostics are computed on finite
historical observations:

The same matured calibration sample is used to estimate IC sign, orient the
signal, calculate the three ranking diagnostics, and select the top five.
Consequently, point-in-time truncation prevents future-data leakage but does
not provide an independent inner/outer assessment of the adaptive selection
rule. The equations below define the implemented selector; they should not be
read as post-selection-unbiased estimates of an individual factor's skill.

\paragraph{Information Coefficient (IC).}
Following Grinold and Kahn \cite{grinold1999active}, the Spearman rank correlation between the normalised alpha signal and
forward returns over the evaluation suffix:
\begin{equation}
    \mathrm{IC}_k \;=\;
    \operatorname{SpearmanCorr}\!\bigl(
        \tilde{\alpha}^{(k)}_{\mathcal{T}_{\mathrm{eval}}},\;
        r^{\mathrm{net},(L)}_{\mathcal{T}_{\mathrm{eval}}}
    \bigr).
    \label{eq:ic}
\end{equation}

\paragraph{Directional Accuracy.}
The fraction of bars where the predicted direction agrees with the
realised direction, restricted to bars with a strong enough signal:
\begin{equation}
    \mathrm{Acc}_k \;=\;
    \frac{1}{|\mathcal{T}^*|}
    \sum_{u \in \mathcal{T}^*}
    \mathbf{1}\!\left[
        \operatorname{sign}\!\bigl(\tilde{\alpha}^{(k)}_u\bigr)
        = y^{(L)}_u
    \right],
\end{equation}
where $\mathcal{T}^* = \bigl\{u : |\tilde{\alpha}^{(k)}_u| > 0.05\bigr\}$
excludes near-zero signals and $y^{(L)}_u$ is the canonical binary label from
Eq.~\eqref{eq:label}; hence a zero net return is \texttt{DOWN}, not a third class.

\paragraph{Long-Only Win Rate.}
The precision of the alpha on bullish signals:
\begin{equation}
    \mathrm{LongAcc}_k \;=\;
    \frac{
        \sum_u \mathbf{1}\!\left[
            \tilde{\alpha}^{(k)}_u > 0.05 \;\wedge\; r^{\mathrm{net},(L)}_u > 0
        \right]
    }{
        \sum_u \mathbf{1}\!\left[
            \tilde{\alpha}^{(k)}_u > 0.05
        \right]
    }.
\end{equation}

\paragraph{Horizon Sharpe diagnostic.}
For ranking, define $z^{(k)}_u=\operatorname{sign}(\tilde{\alpha}^{(k)}_u)
r^{\mathrm{net},(L)}_u$ on the finite historical sample. Following Sharpe
\cite{sharpe1966mutual}, the diagnostic is:
\begin{equation}
    \mathrm{Sharpe}_k \;=\;
    \sqrt{\frac{252}{L}} \cdot
    \frac{
        \mathbb{E}_{u\in\mathcal{T}_{\mathrm{eval}}}\!\left[z^{(k)}_u\right]
    }{
        \mathrm{Std}_{u\in\mathcal{T}_{\mathrm{eval}}}\!\left[z^{(k)}_u\right]
        + \varepsilon
    }.
    \label{eq:sharpe}
\end{equation}
Sortino, maximum drawdown, hit rate, long-only accuracy, and average trade are
retained as descriptive candidate diagnostics but do not enter the composite
selection score.

Candidates with no fitted transform, fewer than 20 finite evaluation
observations, or a non-finite value for any ranking metric are excluded before scaling. Let
$\mathcal{V}\subseteq\mathcal{A}$ be the remaining eligible set and let
$q\in\{|\mathrm{IC}|,\mathrm{Acc},\mathrm{Sharpe}\}$.
The implementation applies candidate-wise min--max scaling:
\begin{equation}
    \mathcal{N}_{\mathcal{V}}(q_k) =
    \begin{cases}
        \dfrac{q_k-q_{\min}}{q_{\max}-q_{\min}},
        & q_{\max}-q_{\min} \geq \varepsilon_{\mathrm{rank}},\\[6pt]
        \dfrac{1}{2}, & \text{otherwise},
    \end{cases}
    \label{eq:candidate_minmax}
\end{equation}
where $q_{\min}=\min_{j\in\mathcal{V}}q_j$,
$q_{\max}=\max_{j\in\mathcal{V}}q_j$, and
$\varepsilon_{\mathrm{rank}}=10^{-9}$. The default composite selection
score is therefore:
\begin{equation}
    S_k \;=\;
    0.40\,\mathcal{N}_{\mathcal{V}}(|\mathrm{IC}_k|)
    \;+\; 0.35\,\mathcal{N}_{\mathcal{V}}(\mathrm{Acc}_k)
    \;+\; 0.25\,\mathcal{N}_{\mathcal{V}}(\mathrm{Sharpe}_k).
    \label{eq:composite}
\end{equation}

The weights prioritise association and directional accuracy while retaining a
risk-adjusted return diagnostic. Long-only accuracy is reported but excluded
from ranking to avoid favouring factors solely because the sample trends up.

Signal orientation is inferred from the sign of IC on the cutoff-safe
historical snapshot. A negative-IC candidate is multiplied by $-1$ for its
historical directional metrics and for the current inference value. All labels
used for this operation have matured by the decision cutoff, but estimating
orientation and ranking on the same finite history remains a selection risk
discussed in Section~\ref{sec:limitations_futurework}.

Let $K=\min(5,|\mathcal{V}|)$. The top candidates are selected as:
\begin{equation}
    \mathcal{A}^{*} \;=\;
    \underset{\mathcal{B}\,\subseteq\,\mathcal{V},\;|\mathcal{B}|=K}
    {\arg\max}
    \;\sum_{\alpha^{(k)} \in \mathcal{B}} S_k.
    \label{eq:selection}
\end{equation}

\subsection{Sentiment Signal Construction}

Let $\mathcal{D}_e = \{d_1, d_2, \ldots, d_{|\mathcal{D}_e|}\}$ be the
set of CafeF articles available to the system for asset $s$ at the timestamp
$\tau_d$ of decision bar $d=e-1$.
Each article is assigned a score by the scorer actually recorded for that
article:
\begin{equation}
    \rho_i \;=\; \mathrm{Scorer}_i(d_i) \;\in\; [-1,\, 1],
\end{equation}
where $\rho_i = +1$ denotes strongly positive, $-1$ strongly negative,
and $0$ neutral sentiment. For newly scored records,
$\mathrm{Scorer}_i$ is either the configured
\texttt{5CD-AI/Vietnamese-Sentiment-visobert} HF Inference API endpoint or
the versioned deterministic lexicon fallback. Each record stores the scorer,
requested model identifier, observed Hub revision and its verification status,
and fallback status/reason. Legacy scores without these fields are marked
\texttt{legacy-unknown}; no model identity is inferred from their numeric values.

The unweighted mean sentiment score for descriptive reporting is:
\begin{equation}
    \bar{\rho}_e \;=\;
    \frac{1}{|\mathcal{D}_e|} \sum_{i=1}^{|\mathcal{D}_e|} \rho_i.
\end{equation}

For the causal sentiment context, each dated article receives an exponential
weight based on its age at cutoff $\tau_d$:
\begin{align}
    \Delta_i(\tau_d) &= \tau_d - \mathrm{date}(d_i), \\
    w_i(\tau_d) &= 2^{-\Delta_i(\tau_d)/h}, \qquad h=30\ \text{days}, \\
    S_{\mathrm{decay},e}
    &= \frac{\sum_{d_i\in\mathcal{D}_e} w_i(\tau_d)\rho_i}
             {\sum_{d_i\in\mathcal{D}_e} w_i(\tau_d)},
    \label{eq:sent_decay} \\[4pt]
    \mathrm{Rel}_{\mathrm{sent},e}
    &\;=\;
    S_{\mathrm{decay},e} \;-\;
    \frac{1}{|\mathcal{R}_e|}
    \sum_{s' \in \mathcal{R}_e} S^{(s')}_{\mathrm{decay},e},
    \label{eq:relsent}
\end{align}
where $\mathcal{R}_e$ contains only related companies with at least three
valid dated-and-scored articles at cutoff $\tau_d$.  The exponential mean makes recent articles
more influential without comparing a batch mean against the same batch from
which it was calculated.  $\mathrm{Rel}_{\mathrm{sent},e}$ compares like-for-like
time-decayed scores.

\paragraph{Batch definition and low-news stability.}
The set $\mathcal{D}_e$ contains all dated articles for target stock $s$ in the 90-day sliding window ending at cutoff $\tau_d$. The complete compact observation set is used for Eq.~\ref{eq:sent_decay}; the 15-article list stored for display is not the normalization sample. With fewer than three valid dated-and-scored observations, $S_{\mathrm{decay},e}$ is marked unavailable rather than assigned a neutral zero. Likewise, $\mathrm{Rel}_{\mathrm{sent},e}$ is unavailable when $\mathcal{R}_e$ is empty. The effective sample size, invalid-date/score counts, and number of post-cutoff observations excluded are stored in run provenance. Neither context variable enters any numerical alpha formula or the registry ranking.

To prevent look-ahead leakage during backtesting, every cached article stores
timezone-aware \texttt{published\_at} and \texttt{available\_at} fields in
\texttt{Asia/Ho\_Chi\_Minh}. The cache is keyed by $(s,\,\tau_d)$ and returns only
articles satisfying $\mathrm{available\_at}(d_i)\leq \tau_d$. CafeF records having
date-only precision are assigned a conservative availability time of 00:00 on
the following local day; records without a parseable date are rejected in
research mode and never used to fill a sparse historical window. All 180
articles in the nine supplied cache artifacts contain parseable dates, so the
previous undated fallback was a latent risk rather than evidence that those
reported runs actually consumed undated articles.

\subsection{Multi-Agent Decision Function}

The full system decomposes $\mathcal{F}$ into five
sub-functions corresponding to the five specialised agents:
\begin{align}
    \mathbf{r}^{\mathrm{ind}}
    &\;=\; \mathcal{F}_{\mathrm{ind}}\!\bigl(\mathbf{X}_e\bigr),
    \label{eq:ind} \\
    \mathbf{a}_e
    &\;=\; \mathcal{F}_{\mathrm{factor}}\!\bigl(
            \mathbf{X}_e,\; \mathcal{A}^{*}
           \bigr),
    \label{eq:alpha_factor} \\
    \mathbf{c}^{\mathrm{sent}}_e
    &\;=\; \mathcal{F}_{\mathrm{context}}\!\bigl(
            \mathcal{D}_e,\;
            S_{\mathrm{decay},e},\;
            \mathrm{Rel}_{\mathrm{sent},e}
           \bigr),
    \label{eq:sent_context} \\
    \mathbf{r}^{\mathrm{alpha}}
    &\;=\; \mathcal{F}_{\mathrm{text}}\!\bigl(
            \mathbf{a}_e,\; \mathbf{c}^{\mathrm{sent}}_e
           \bigr),
    \label{eq:alpha} \\
    \mathbf{r}^{\mathrm{pat}}
    &\;=\; \mathcal{F}_{\mathrm{pat}}\!\bigl(
            \mathbf{I}^{\mathrm{kline}}_e
           \bigr),
    \label{eq:pat} \\
    \mathbf{r}^{\mathrm{trend}}
    &\;=\; \mathcal{F}_{\mathrm{trend}}\!\bigl(
            \mathbf{I}^{\mathrm{trend}}_e
           \bigr),
    \label{eq:trend}
\end{align}
where $\mathbf{I}^{\mathrm{kline}}_e$ and
$\mathbf{I}^{\mathrm{trend}}_e$ are the candlestick and
trendline-annotated chart images rendered from $\mathbf{X}_e$, and
$\mathbf{c}^{\mathrm{sent}}_e$ is a separate textual context containing the
target and related-company sentiment summaries. The factor values in
$\mathbf{a}_e$ are functions of $\mathbf{X}_e$ only; the LLM compares
them qualitatively with $\mathbf{c}^{\mathrm{sent}}_e$ when producing
$\mathbf{r}^{\mathrm{alpha}}$.

The Decision Agent aggregates all four reports through structured
chain-of-thought prompting \cite{wei2022chain} and produces the final binary signal:
\begin{equation}
    \hat{y}_{e,L} \;=\; \mathcal{F}_{\mathrm{dec}}\!\bigl(
        \mathbf{r}^{\mathrm{ind}},\;
        \mathbf{r}^{\mathrm{alpha}},\;
        \mathbf{r}^{\mathrm{pat}},\;
        \mathbf{r}^{\mathrm{trend}}
    \bigr) \;\in\; \{-1,\, +1\}.
    \label{eq:decision}
\end{equation}

The forced binary output constraint ensures $\hat{y}_{e,L}\neq 0$ for every
valid \emph{prediction}; it does not imply an always-invested account. Under
the default long-only action map, \texttt{SHORT} becomes \texttt{CASH}.
\texttt{LONG} executes a complete next-open-to-target-close cycle when valid
positive prices and timestamps are available. This separation permits binary
directional evaluation without misrepresenting a bearish prediction as an
executed short sale.

\subsection{Walk-Forward Evaluation Protocol}

We adopt a rolling walk-forward evaluation protocol commonly used in quantitative trading system evaluation \cite{pardo2008evaluation}.

Given a historical series of $T$ bars, valid test points are exclusive
boundaries with $W\leq e\leq T-L$.  The implementation uses the latest
available boundaries at spacing $S$ (up to $N_{\mathrm{test}}$ points) and
processes them in chronological order:
\begin{equation}
    \mathcal{T} \;=\;
    \bigl\{\, e_i = T-L-(i-1)\cdot S
    \;\mid\; i = 1, 2, \ldots, N_{\mathrm{test}} \,\bigr\},
    \label{eq:testpoints}
\end{equation}
retaining only values $e_i\geq W$ before reversing the list for execution.
Thus $e=W$ is included whenever it is among the requested points, and every
selected point satisfies the complete-window and complete-target conditions.
Here $S$ is the step size between consecutive test points and
$N_{\mathrm{test}} = 20$ in our main experiments.

The three primary evaluation metrics are:
\begin{align}
    \mathrm{Acc}
    &\;=\;
    \frac{1}{|\mathcal{T}|}
    \sum_{e \in \mathcal{T}}
    \mathbf{1}\!\bigl[\hat{y}_{e,L} = y_{e,L}\bigr],
    \label{eq:acc} \\[4pt]
    \Delta\mathrm{Acc}
    &\;=\;
    \mathrm{Acc}_{\mathrm{Full}}
    \;-\; \mathrm{Acc}_{\mathrm{No\text{-}\alpha}},
    \label{eq:lift} \\[4pt]
    q_e
    &\;=\;
    a_e\left[
    \frac{C_{e-1+L}}{O_e}
    \frac{(1-c_{\mathrm{slip}})(1-c_{\mathrm{fee}})}
         {(1+c_{\mathrm{slip}})(1+c_{\mathrm{fee}})}-1
    \right],
    \label{eq:execution_return}\\[4pt]
    W_{k+1}
    &\;=\;
    W_k\left(1+q_{e_k}\right),\qquad W_0=1,
    \label{eq:wealth_update}\\[4pt]
    R_{\mathrm{acct}}
    &\;=\;
    \frac{W_K}{W_0}-1,
    \label{eq:account_return}
\end{align}
where $\Delta\mathrm{Acc}$ (\emph{Alpha Lift}) quantifies the improvement
in directional accuracy attributable to the Alpha Agent. In the economic
simulation, $a_e=1$ for \texttt{LONG} and $a_e=0$ for \texttt{SHORT}. Each
LONG allocates all available cash at $O_e$, pays fee and slippage on the BUY
leg, liquidates all shares at $C_{e-1+L}$, and pays the same rates on the SELL
leg. The configured values are $c_{\mathrm{fee}}=0.25\%$ and
$c_{\mathrm{slip}}=0.10\%$ per leg. Each test point is a closed cycle and the
implementation requires $S\geq L$, so no shares carry into the next point.
SHORT forecasts hold cash and incur neither fee nor slippage. Let
$\rho_k=W_{k+1}/W_k-1=q_{e_k}$ denote the cycle return, including zeros for
CASH. The account risk metrics are
\begin{align}
    \mathrm{Sharpe}_{\mathrm{acct}}
    &=\sqrt{252}\,\frac{\mathbb{E}[\rho_k]}
    {\mathrm{Std}[\rho_k]+\varepsilon},
    \label{eq:account_sharpe}\\[4pt]
    \mathrm{MDD}_{\mathrm{acct}}
    &=\max_{0\leq k\leq K}
      \left(1-\frac{W_k}{\max_{0\leq i\leq k}W_i}\right).
    \label{eq:account_mdd}
\end{align}
The Sharpe risk-free rate is zero, and the population standard deviation is
used. Maximum drawdown is a relative wealth loss and includes the initial
wealth anchor $W_0$, so an initial losing cycle is not omitted.

Finally, let $\mathcal{X}=\{e_k:a_{e_k}=1\}$ be the executed-position sample.
Trade-level summaries are
\begin{equation}
    \mathrm{HitRate}=\frac{1}{|\mathcal{X}|}
    \sum_{e\in\mathcal{X}}\mathbf{1}[q_e>0],\qquad
    \mathrm{AvgTrade}=\frac{1}{|\mathcal{X}|}
    \sum_{e\in\mathcal{X}}q_e,
    \label{eq:account_trade_metrics}
\end{equation}
with both defined as zero when $\mathcal{X}$ is empty. SHORT/CASH cycles do not
dilute either executed-trade denominator. Risk ratios are reported as zero when
their dispersion denominator does not exceed the numerical tolerance. The
cost rates are research allowances rather than a claim about a particular
venue's realised fee schedule.
